% ==============================================================================
% 05_criteria2_structure.tex — Criteria 2: Programme Structure and Content
% Score: 2
% ==============================================================================

\criterion{2}{โครงสร้างและเนื้อหาของหลักสูตร}{Programme Structure and Content}

% ------------------------------------------------------------------------------
\criterionsub{2.1 The specifications of the programme and all its courses are shown to be comprehensive, up-to-date, and made available and communicated to all stakeholders.}

หลักสูตรวิศวกรรมคอมพิวเตอร์และปัญญาประดิษฐ์ บัณฑิตศึกษา พ.ศ. 2568 มีโครงสร้างหลักสูตรครบถ้วนตามเกณฑ์ มคอ.2 ครอบคลุมหมวด 1--16 รวมรายวิชาบังคับ 6 วิชา และวิทยานิพนธ์/สารนิพนธ์ 5 วิชา สื่อ/เอกสารที่ใช้เผยแพร่และสื่อสารหลักสูตรไปยังผู้มีส่วนได้ส่วนเสียประกอบด้วย
โดยหลักสูตรใช้ช่องทางสื่อสารหลายรูปแบบ ได้แก่ \textbf{มคอ.2 ฉบับผ่านสภา}ที่เผยแพร่บนเว็บไซต์มหาวิทยาลัยเพื่อให้ผู้ที่สนใจและผู้สมัครเข้าศึกษาเข้าถึงได้ \textbf{แผ่นพับประชาสัมพันธ์}สำหรับงานเปิดบ้าน/งานศิษย์เก่าซึ่งสรุป PLO โครงสร้างหลักสูตร และเส้นทางอาชีพ \textbf{เว็บไซต์คณะเทคโนโลยีอุตสาหกรรม}ที่แสดงหน้าหลักสูตร CE\&AI ซึ่งมีการอัปเดตข้อมูล และ \textbf{การปฐมนิเทศ}ที่ใช้ชี้แจง PLO/CLO และแผนการศึกษาให้กับนักศึกษาใหม่อย่างเป็นทางการ

ข้อกำหนดของรายวิชา (Course Specification, มคอ.3) ได้ครอบคลุม CLO กิจกรรมการเรียนรู้ และการประเมินผล โดยอยู่ระหว่างการจัดทำให้ครบทุกรายวิชาก่อนเปิดสอน

\marknote{มคอ.3 ฉบับเต็มของทุกรายวิชายังไม่แล้วเสร็จ ต้องดำเนินการให้ครบก่อนเปิดสอนภาคแรก}

% ------------------------------------------------------------------------------
\criterionsub{2.2 The design of the curriculum is shown to be constructively aligned with achieving the expected learning outcomes}

หลักสูตรออกแบบ Constructive Alignment ระหว่าง CLO $\rightarrow$ TLO $\rightarrow$ PLO อย่างเป็นระบบ แต่ละรายวิชามีการกำหนดกิจกรรมการเรียนรู้และวิธีการประเมินที่สอดคล้องกับ CLO และ CLO ทุกตัวสามารถสืบไปถึง PLO ที่เกี่ยวข้องได้

\begin{table}[H]
  \centering
  \caption{ตัวอย่าง Constructive Alignment ของวิชา 4095101 (AI \& ML)}
  \label{tbl:clo_plo_example}
  \begin{tabularx}{\linewidth}{|c|>{\raggedright\arraybackslash}p{2.4cm}|c|c|Y|}
    \hline
    \rowcolor{gray!15}
    \textbf{CLO} & \textbf{Action verb} & \textbf{Bloom} & \textbf{PLO} & \textbf{การประเมิน} \\
    \hline
    CLO1 & อธิบาย       & U  & PLO1 & สอบกลางภาค \\
    CLO2 & ประยุกต์ใช้ & Ap & PLO1 & ปฏิบัติการ \\
    CLO3 & วิเคราะห์   & An & PLO3 & รายงานกลุ่ม \\
    CLO4 & ประเมินค่า  & Va & PLO3 & สัมมนา \\
    CLO5 & พัฒนา       & A  & PLO1, PLO3 & โปรเจกต์ \\
    \hline
  \end{tabularx}
\end{table}

% ------------------------------------------------------------------------------
\criterionsub{2.3 The design of the curriculum is shown to include feedback from stakeholders, especially external stakeholders.}

การออกแบบหลักสูตรมีการรับฟังข้อเสนอแนะจาก (1) สภาวิชาการมหาวิทยาลัย (2) คณะกรรมการประจำคณะเทคโนโลยีอุตสาหกรรม (3) สภามหาวิทยาลัยราชภัฏอุตรดิตถ์ (มติครั้งที่ 4/2568) และ (4) ผู้ทรงคุณวุฒิภายนอกที่ร่วมวิพากษ์หลักสูตร อย่างไรก็ตาม การรับฟังความคิดเห็นจาก external stakeholder ภาคอุตสาหกรรม/ผู้ประกอบการในรูปแบบการประชุมเชิงปฏิบัติการอย่างเป็นทางการยังมีหลักฐานไม่ครบถ้วน

\marknote{ต้องการหลักฐาน: รายงานสรุปการประชุมรับฟังความคิดเห็นจากภาคอุตสาหกรรม/ผู้ใช้บัณฑิต ขณะยกร่างหลักสูตร}

% ------------------------------------------------------------------------------
\criterionsub{2.4 The contribution made by each course in achieving the expected learning outcomes is shown to be clear.}

ตาราง CLO--PLO Mapping แสดงความเชื่อมโยงของแต่ละรายวิชากับ PLO อย่างชัดเจน

\begin{landscape}
\begin{table}[H]
  \centering
  \caption{ตาราง CLO--PLO Mapping วิชาบังคับและวิทยานิพนธ์/สารนิพนธ์}
  \label{tbl:clo_plo_map}
  \begin{tabularx}{\linewidth}{|>{\raggedright\arraybackslash}X|Y|Y|Y|}
    \hline
    \rowcolor{gray!15}
    \textbf{รหัสวิชา} & \textbf{ชื่อวิชา} & \textbf{Bloom's Levels} & \textbf{PLOs} \\
    \hline
    4095101 & ปัญญาประดิษฐ์และการเรียนรู้ของเครื่อง & U, Ap, An, Va, A & PLO1, PLO3 \\
    4095102 & การเขียนโปรแกรมขั้นสูงสำหรับ ML       & U, Ap, An, Va, A & PLO1, PLO3 \\
    7015906 & ระเบียบวิธีวิจัยทางวิทย์/วิศวกรรม      & U, An, O, In, A  & PLO2, PLO5 \\
    7015907 & สัมมนาวิศวกรรม CE-AI                   & An, E, O, In, P  & PLO2, PLO4, PLO5 \\
    4095101 & เทคโนโลยีอุบัติใหม่ทาง CE\&AI          & U, Ap, Re, Rs, P & PLO1, PLO3 \\
    7015908 & หัวข้อพิเศษทาง CE\&AI                  & U, An, E, Va, OA & PLO1, PLO3, PLO4 \\
    \hline
    7015901 & วิทยานิพนธ์ 1 (แผน 1) & An, Cr, O, In, A  & PLO1, PLO3, PLO4, PLO5 \\
    7015902 & วิทยานิพนธ์ 2 (แผน 1) & An, Cr, In, A     & PLO1, PLO3, PLO5 \\
    7015903 & วิทยานิพนธ์ 3 (แผน 1) & E, Cr, In, A      & PLO3, PLO4, PLO5 \\
    7015904 & สารนิพนธ์ 1 (แผน 2)   & Ap, An, OA        & PLO1, PLO2, PLO3, PLO5 \\
    7015905 & สารนิพนธ์ 2 (แผน 2)   & Ap, E, In, A      & PLO1, PLO3, PLO4, PLO5 \\
    \hline
  \end{tabularx}
\end{table}
\end{landscape}

% ------------------------------------------------------------------------------
\criterionsub{2.5 The curriculum to show that all its courses are logically structured, properly sequenced (progression from basic to intermediate to specialised courses), and are integrated.}

ลำดับวิชาในหลักสูตรเรียงจากพื้นฐาน $\rightarrow$ ขั้นกลาง $\rightarrow$ ขั้นสูง/บูรณาการ

โดยในช่วง \textbf{Basic} (ภาคต้น ปี 1) ผู้เรียนเริ่มจากรายวิชาพื้นฐานสำคัญ ได้แก่ 4095101 ปัญญาประดิษฐ์ฯ 4095102 การเขียนโปรแกรมขั้นสูง และ 7015906 ระเบียบวิธีวิจัยฯ จากนั้นเข้าสู่ช่วง \textbf{Intermediate} (ภาคปลาย ปี 1) ผ่าน 7015907 สัมมนา และ 4095101 เทคโนโลยีอุบัติใหม่ฯ ก่อนพัฒนาไปสู่ช่วง \textbf{Specialised/Integration} (ปี 2) ซึ่งเน้นการบูรณาการความรู้ใน 7015908 หัวข้อพิเศษ และวิทยานิพนธ์/สารนิพนธ์

วิชาบูรณาการ (integrated) ที่นำความรู้จากหลายวิชาและศาสตร์อื่นมาใช้ร่วม ได้แก่ 7015908 หัวข้อพิเศษฯ และวิทยานิพนธ์/สารนิพนธ์ ซึ่งใช้ทั้งความรู้ AI/ML ระเบียบวิธีวิจัย สถิติ และจริยธรรมวิชาชีพ

% ------------------------------------------------------------------------------
\criterionsub{2.6 The curriculum to have option(s) for students to pursue major and/or minor specialisations.}

หลักสูตรมี 2 แผนชัดเจนเปิดโอกาสให้ผู้เรียนเลือกตามความสนใจ
โดย \textbf{แผน 1 (วิชาการ)} เน้นวิทยานิพนธ์ 3 ภาค (7015901--3) เหมาะกับผู้ที่ต้องการสร้างองค์ความรู้ใหม่และมีเป้าหมายศึกษาต่อระดับปริญญาเอก ขณะที่ \textbf{แผน 2 (วิชาชีพ)} เน้นสารนิพนธ์ 2 ภาค (7015904--5) เหมาะกับผู้ที่ต้องการพัฒนาต้นแบบนวัตกรรมและนำไปใช้ในงานจริงในภาคอุตสาหกรรม

ทั้งสองแผนมีหน่วยกิตรวมไม่น้อยกว่า 36 หน่วยกิต โดยโครงสร้างมคอ.2 สรุปได้ว่าในแผน 1 เป็นรายวิชาเฉพาะด้านบังคับ 18 หน่วยกิต วิชาเลือก 6 หน่วยกิต และวิทยานิพนธ์ 12 หน่วยกิต ส่วนแผน 2 เป็นรายวิชาเฉพาะด้านบังคับ 18 หน่วยกิต วิชาเลือก 12 หน่วยกิต และสารนิพนธ์ 6 หน่วยกิต (มีวิชาเสริมไม่นับหน่วยกิต 6 หน่วยกิตเหมือนกัน) ซึ่งทำให้ผู้เรียนเลือกเส้นทางการเรียนรู้ได้ตามเป้าหมายอาชีพและการพัฒนาตนเอง (อ้างอิง AUNQA-2-1)

นอกจากเลือกแผนแล้ว ผู้เรียนยังเลือกอาจารย์ที่ปรึกษาวิทยานิพนธ์/สารนิพนธ์ตามหัวข้อความสนใจ และเลือกหัวข้อวิจัยได้อย่างอิสระ ภายใต้การกำกับของคณะกรรมการบัณฑิตศึกษา

\begin{table}[H]
  \centering
  \caption{ตารางเปรียบเทียบโครงสร้างหลักสูตร แผน 1 vs แผน 2}
  \begin{tabularx}{\linewidth}{|>{\raggedright\arraybackslash}X|c|c|}
    \hline
    \textbf{องค์ประกอบ} & \textbf{แผน 1 (วิชาการ)} & \textbf{แผน 2 (วิชาชีพ)} \\
    \hline
    วิชาบังคับ                     & 6 วิชา & 6 วิชา \\
    วิทยานิพนธ์/สารนิพนธ์           & 3 ภาค (36 นก.) & 2 ภาค \\
    \textbf{รวม}                   & \textbf{36 นก.} & \textbf{36 นก.} \\
    \hline
  \end{tabularx}
\end{table}

% ------------------------------------------------------------------------------
\criterionsub{2.7 The programme to show that its curriculum is reviewed periodically following an established procedure and that it remains up-to-date and relevant to industry.}

หลักสูตรกำหนดงวดการประเมินและทบทวนหลักสูตรทุก 5 ปี ตามข้อบังคับมหาวิทยาลัยราชภัฏอุตรดิตถ์ ว่าด้วยการจัดการศึกษาระดับบัณฑิตศึกษา พ.ศ. 2566 และอาจทบทวนรอบเล็กทุก 1 ปีการศึกษาเพื่อปรับเนื้อหาวิชาให้ทันสมัย วิธีการที่ใช้ในการประเมินและทบทวนประกอบด้วย
การประเมินรายวิชาและการสอนโดยนักศึกษาทุกภาคการศึกษา การประชุมทบทวน CLO/PLO ประจำปี และการสำรวจความพึงพอใจของผู้ใช้บัณฑิต/ศิษย์เก่าเมื่อมีบัณฑิตแล้ว เพื่อให้ข้อมูลย้อนกลับเพียงพอสำหรับการปรับปรุงอย่างเป็นระบบ

อย่างไรก็ตาม เนื่องจากหลักสูตรนี้เพิ่งผ่านสภาในปี 2568 จึงยังไม่มีผลการทบทวนหลักสูตรจริง

\begin{selfassessmentbox}{2}{หลักสูตรมีการออกแบบโครงสร้างที่ครบถ้วน CLO--PLO Mapping ชัดเจน ลำดับวิชาเหมาะสม มีตัวเลือก 2 แผน และมีระบบทบทวนหลักสูตรตามข้อบังคับ แต่ มคอ.3 ฉบับเต็มยังไม่ครบ ขาดหลักฐาน external stakeholder feedback อย่างเป็นทางการ และยังไม่มีผลการทบทวนหลักสูตรจริง (หลักสูตรใหม่)}
\end{selfassessmentbox}

% ------------------------------------------------------------------------------
\subsection*{รายการหลักฐาน}

หลักฐานประกอบ Criterion~2 ประกอบด้วย AUNQA-2-1 มคอ.2 ฉบับผ่านสภามหาวิทยาลัยราชภัฏอุตรดิตถ์, AUNQA-2-2 รวมเล่ม CLO รายวิชา พร้อม CLO--PLO Mapping, AUNQA-2-3 ข้อบังคับมหาวิทยาลัยราชภัฏอุตรดิตถ์ ว่าด้วยการจัดการศึกษาระดับบัณฑิตศึกษา พ.ศ. 2566 และ AUNQA-2-4 มคอ.3 ของรายวิชาบังคับ (อยู่ระหว่างจัดทำ)
